\label{chap:planck-autodual}

\section{Realización de Compton asociada a la periodicidad temporal de \(108\) fases}

Sobre el ciclo de \cref{eq:realizacion-circular-ct108}, reutilizamos la sección
positiva de acción \(\hbar\) ya generada por la rama determinantal de acción.
El cuanto angular asociado es
\begin{equation}
E_{108}=\hbar\omega_{108},
\qquad
m_{108}=\frac{E_{108}}{c^2}
=\frac{\hbar\omega_{108}}{c^2}.
\label{eq:cuanto-ct108}
\end{equation}
No se ha empleado todavía \(G\). La geometría del ciclo y la cuantización
angular bastan para fijar sus dos lectores de Compton.

\begin{theorem}[Identidad entre la periodicidad temporal de \(108\) fases y la longitud de Compton]
\label{thm:ct108-compton}
Para el cuanto \eqref{eq:cuanto-ct108},
\begin{equation}
\lambdabarC(m_{108})=\frac{\hbar}{m_{108}c}=R_{108},
\qquad
\lambdaC(m_{108})=\frac{h}{m_{108}c}=C_{108}.
\label{eq:identidad-ct108-compton}
\end{equation}
Por tanto, el radio del ciclo es su longitud de Compton reducida y la
longitud de una vuelta es su longitud de Compton completa.
\end{theorem}

\begin{proof}
Sustituyendo \(m_{108}=\hbar\omega_{108}/c^2\),
\[
\frac{\hbar}{m_{108}c}
=\frac{\hbar c^2}{\hbar\omega_{108}c}
=\frac{c}{\omega_{108}}=R_{108}.
\]
Como \(h=2\pi\hbar\), la segunda igualdad es \(2\pi R_{108}=C_{108}\).
\end{proof}

La coincidencia no es un ajuste numérico. Es una identidad homogénea: vale
para todo \(\hbar,c,t_0>0\) y depende sólo de que el mismo periodo determine
la frecuencia, la longitud del ciclo y el cuanto de acción.

\section{Transductor gravitatorio y ecuación de selección autodual}

La recta gravitatoria HMT--MD contiene el transductor
\begin{equation}
G(\theta)=\theta\frac{c^5t_0^2}{\hbar},
\qquad \theta>0.
\label{eq:transductor-g-planck-anexo}
\end{equation}
La homogeneidad es exacta, pero no selecciona por sí sola \(\theta\). El
ciclo CT108 aporta una condición interna adicional: el lector gravitatorio
debe devolver la misma escala radial que el lector geométrico y el lector de
Compton.

Definimos el radio gravitatorio reducido
\begin{equation}
\rg(\theta):=\frac{G(\theta)m_{108}}{c^2}.
\label{eq:radio-gravitatorio-reducido}
\end{equation}
Al sustituir \eqref{eq:cuanto-ct108} y
\eqref{eq:transductor-g-planck-anexo},
\begin{equation}
\rg(\theta)
=\theta\frac{\pi}{54}\ell_0.
\label{eq:rg-theta-lineal}
\end{equation}

\begin{theorem}[Selector autodual de la periodicidad temporal de \(108\) fases]
\label{thm:selector-autodual-ct108}
La ecuación
\[
\rg(\theta)=R_{108}
\]
posee una única solución positiva:
\begin{equation}
\boxed{\thetaStar=\left(\frac{54}{\pi}\right)^2.}
\label{eq:theta-star}
\end{equation}
\end{theorem}

\begin{proof}
Por \eqref{eq:realizacion-circular-ct108} y
\eqref{eq:rg-theta-lineal}, la ecuación es
\[
\theta\frac{\pi}{54}\ell_0
=\frac{54}{\pi}\ell_0.
\]
Como \(\ell_0>0\), se simplifica y se obtiene
\(\theta=(54/\pi)^2\). La función
\(\theta\mapsto\rg(\theta)\) es estrictamente creciente en
\(\RR_{>0}\), de modo que no existe otra solución positiva.
\end{proof}

\begin{corollary}[Unidad de Planck interna]
\label{cor:unidad-planck-interna}
Con \(G_*=G(\thetaStar)\), las secciones de Planck satisfacen
\begin{equation}
\ellP:=\sqrt{\frac{\hbar G_*}{c^3}}=R_{108},
\qquad
\tP:=\sqrt{\frac{\hbar G_*}{c^5}}=\frac1{\omega_{108}},
\label{eq:planck-long-time-ct108}
\end{equation}
\begin{equation}
\mP:=\sqrt{\frac{\hbar c}{G_*}}=m_{108},
\qquad
\EP:=\mP c^2=E_{108}.
\label{eq:planck-mass-energy-ct108}
\end{equation}
En particular,
\begin{equation}
\boxed{
R_{108}=\lambdabarC=\rg=\ellP,
\qquad
C_{108}=\lambdaC=2\pi\rg=2\pi\ellP.}
\label{eq:tres-lectores-planck}
\end{equation}
\end{corollary}

\begin{proof}
Se sustituye \(G_*=(54/\pi)^2c^5t_0^2/\hbar\). Entonces
\[
\ellP=\frac{54}{\pi}ct_0=R_{108},
\qquad
\tP=\frac{54}{\pi}t_0=\frac1{\omega_{108}},
\]
y
\[
\mP=\frac{\pi}{54}\frac{\hbar}{c^2t_0}
=\frac{\hbar\omega_{108}}{c^2}=m_{108}.
\]
El resto se sigue del \cref{thm:ct108-compton}.
\end{proof}

\begin{remark}[Tres lectores, una sección]
La geometría circular publica \(R_{108}\) y \(C_{108}\); la cuantización de
Compton publica \(\lambdabarC\) y \(\lambdaC\); la realización gravitatoria
publica \(\rg\) y \(2\pi\rg\). En \(\thetaStar\) esas tres parejas son
coordenadas de una única sección. Esta es la razón precisa por la que la
unidad de Planck deja de aparecer como una combinación dimensional accidental
dentro de esta formalización.
\end{remark}

\section{Equivalencia de normalizaciones sobre la identidad autodual}

El transductor \eqref{eq:transductor-g-planck-anexo} admite dos formas de
fijar la escala elemental, que no deben confundirse.

\begin{enumerate}[label=\textup{(\roman*)}]
\item Si se fija \(\theta=1\), entonces
\(\sqrt{\hbar G/c^3}=\ell_0\). La transición elemental \(ct_0\) se identifica
con la longitud de Planck de esa carta.
\item Si se exige que la \emph{órbita completa} CT108 posea radio de Planck,
la condición selecciona \(\thetaStar=(54/\pi)^2\) y
\(R_{108}=\ellP\).
\end{enumerate}

No son valores rivales de \(G\). Son dos calibres de la misma ecuación
homogénea: uno asigna la unidad de Planck al paso elemental; el otro, al radio
del ciclo completo. El anexo utiliza el segundo porque su tesis concierne a
la identidad de la clausura CT108.

\section{Convenciones gravitatorias: radio reducido y radio de Schwarzschild}

La notación gravitatoria se fija de manera inequívoca:
\[
\rg=\frac{Gm}{c^2},
\qquad
\rs=\frac{2Gm}{c^2}=2\rg.
\]
El teorema anterior emplea \(\rg\), no \(\rs\). En la convención habitual de
masa de Planck, la igualdad \(\hbar/(mc)=Gm/c^2\) se produce en \(m=\mP\),
mientras
\begin{equation}
\frac{\hbar}{mc}=\frac{2Gm}{c^2}
\quad\Longleftrightarrow\quad
m=\frac{\mP}{\sqrt2}.
\label{eq:cruce-schwarzschild-factor-dos}
\end{equation}
Por ello la frase «escala gravitatoria» designa en la tesis rectora el radio
reducido. El radio de horizonte convencional permanece como lector distinto.

\section{Desdoblamiento bidireccional de la recta de acción}

La recta interna de acción posee dos secciones orientadas:
\begin{equation}
\hbarpre=H_5\,10^{-34}\mathcal U_S,
\qquad
\hbarret=\eta_{\mathrm{ret}}\,10^{-34}\mathcal U_S,
\qquad
\Ract=\frac{H_5}{\eta_{\mathrm{ret}}}.
\label{eq:hbar-pre-ret-anexo}
\end{equation}
La década \(10^{-34}\) fija la escala absoluta; el cociente \(\Ract\)
transporta la diferencia orientada. Esta separación permite preguntar si la
autodualidad de Planck sobrevive a ambas coordenadas.

\begin{theorem}[Covariancia recíproca de acción y gravedad]
\label{thm:accion-gravedad-reciproca}
Para \(a\in\{\mathrm{pre},\mathrm{ret}\}\), sea
\[
G_a(\theta)=\theta\frac{c^5t_0^2}{\hbar_a},
\qquad
m_{108,a}=\frac{\hbar_a\omega_{108}}{c^2}.
\]
Entonces
\begin{equation}
\frac{G_a(\theta)m_{108,a}}{c^2}
=\theta\frac{\pi}{54}\ell_0
\label{eq:invariante-pre-ret-planck}
\end{equation}
es independiente de \(a\). Además,
\begin{equation}
\frac{m_{108,\mathrm{pre}}}{m_{108,\mathrm{ret}}}=\Ract,
\qquad
\frac{G_{\mathrm{pre}}}{G_{\mathrm{ret}}}=\Ract^{-1},
\label{eq:reciprocidad-pre-ret}
\end{equation}
y la longitud de Planck
\(\sqrt{\hbar_aG_a/c^3}=\sqrt\theta\,\ell_0\) es la misma en ambas
orientaciones.
\end{theorem}

\begin{proof}
Los factores \(\hbar_a\) se cancelan en el producto \(G_am_{108,a}\).
Las dos razones de \eqref{eq:reciprocidad-pre-ret} se obtienen directamente
de las definiciones. La invariancia de la longitud de Planck se sigue de
\(\hbar_aG_a=\theta c^5t_0^2\).
\end{proof}

Una parametrización simétrica hace visible la geometría:
\begin{equation}
\hbar_{\mathrm{geom}}:=\sqrt{\hbarpre\hbarret},
\qquad
\xi:=\frac12\log\Ract,
\qquad
\hbarpre=\hbar_{\mathrm{geom}}e^\xi,
\quad
\hbarret=\hbar_{\mathrm{geom}}e^{-\xi}.
\label{eq:rapidez-accion-pre-ret}
\end{equation}
La memoria bidireccional se aloja en \(\xi\); la geometría de Planck se
aloja en el producto invariante. Esto permite una diferencia real entre las
dos coordenadas de acción sin exigir una diferencia de velocidad de
propagación.

\begin{remark}[Alcance físico del teorema]
La reciprocidad de \cref{thm:accion-gravedad-reciproca} es exacta dentro del
transductor. Para identificar \(G_{\mathrm{pre}}\) y \(G_{\mathrm{ret}}\) con dos
mediciones gravitatorias orientadas hace falta especificar el protocolo, el
observable y la simetría que relaciona ambos. El teorema prueba la
compatibilidad estructural; no anticipa por sí solo una violación de Lorentz
ni una variación observable de \(G\).
\end{remark}
